\subsection{Expanded natural-conversation study}
A broader inventory inspected 640 local session files. After removing machine
context, generated subagent messages, detected secret patterns and exact
normalized duplicates, it retained 2,814 explicit user messages from 119 sessions
and 35 project groups: 1,336 messages from Codex and 1,478 from Claude. These
messages contained 2,111,697 characters. All eligible messages were submitted
for provisional annotation; long messages were split explicitly rather than
truncated. Annotation failures were retried in smaller batches. One retry batch
still failed and supplied no accepted labels.

The teacher produced 3,235 candidate spans. Validation excluded 192 invalid
annotations, ten exact candidate duplicates across partitions, and one duplicate
identifier. The remaining 3,032 candidates all fit the complete encoder input
budget. They contained 1,763 temporary, 535 durable, 443 ignored and 291 uncertain
labels. These are policy annotations by a model, with no human-reviewed labels.

A fixed hash ordering assigned 21 whole project groups to training, seven to
development and seven to testing, before classifier fitting. This produced
1,629, 1,040 and 363 natural candidates respectively. Synthetic training and
development partitions were added as in the earlier study. The following table
uses only the natural test partition; 310 candidates have a type label.

\begin{table}[htbp]
\centering
\begin{tabular}{lrrrr}
\toprule
Training & Retention & Type & Accepted & Errors\\
\midrule
Synthetic only, frozen & 73/363 & 186/310 & 108/363 & 78/108\\
Mixed, frozen & 215/363 & 279/310 & 61/363 & 18/61\\
Mixed, adapted & 233/363 & 285/310 & 45/363 & 13/45\\
\bottomrule
\end{tabular}
\caption{Agreement with provisional labels on seven held-out project groups.
The joint decision rule uses the fixed 0.8 thresholds. An accepted error is an
incorrect retention label or an incorrect type for a retained candidate.}
\label{tab:expanded-classification}
\end{table}

Natural examples substantially improved agreement relative to synthetic-only
training, but accepted-prediction disagreement remained approximately 29\% for
both mixed models. The adapted model accepted only 12.4\% of test candidates.
This is insufficient evidence for autonomous retention decisions. The apparent
zero-error result among seven accepted predictions in the smaller study did not
persist at this scale. A live Go API check reproduced offline predictions for
361 candidates with matching context, with a maximum probability difference
below $10^{-6}$ and no memory writes. Two candidates at an explicitly split
source boundary had different complete context and were excluded from this
numerical-equivalence check.

\subsection{Operational and answer evaluation}
The full natural corpus indexed into 7,184 source windows. An observed
whitespace-only window exposed a client-side rejection that had interrupted
indexing; accepting nonempty whitespace preserved complete coverage. The final
local indexing run completed without silently truncating an input.

Large-context consolidation also exposed invalid model-generated source
references. Validation rejected those proposals and retained completed batches.
The final context uses request-local aliases and omits historical source lists.
Aliases resolve to exact IDs before the existing transaction checks. The
operational results below distinguish successful acknowledgement from semantic
correctness of the retained memories.

The recovered full-corpus run acknowledged all 2,814 observations exactly
once, with no missing or foreign source references. A fresh run using local
aliases processed 2,600 observations before a summary-length rejection; a retry
processed the remaining 214. Across the fresh run and retry, fifteen successful
chunks acknowledged all observations, and one rejected chunk acknowledged none.
The original recovery run contained five rejected chunks. These tests establish
observed persistence behavior, not the factual accuracy of every retained memory.

The twelve-case development set was evaluated twice with the frozen reader.
Source review found 11 reference-aligned answers in the first run and 12 in the
second, including respectively three and four explicitly conditional answers.
The first run abstained on an ambiguous museum-visit date. In the second, it
provided the reference value under a stated interpretation of that date. Thus
12/12 was attained in one run; the combined result was 23/24, not evidence of
perfect repeatability.

\subsection{Held-out public evaluation}
Sixty additional LongMemEval cases were selected before implementation changes:
ten per question type, ordered by a fixed SHA256 rule, excluding all nineteen
cases used in earlier development. Full supplied histories were ingested under
separate owner scopes. Questions and reference answers were excluded from
memory ingestion. The reader implementation was frozen before the final
measurement. Earlier partial executions were retained privately and were not
used to select prompts or tune against reference answers.

Both conditions received at most eighty excerpts and 32,000 serialized UTF-8
evidence bytes. The baseline used single-query retrieval and one direct reading
call. The source-checked reader used query planning, extraction and composition,
with at most one additional format repair. This is an equal evidence-size
comparison, not an equal-call-budget comparison. Scoring combined model grading
with Codex review of references and source evidence. No independent human judge
participated. Table~\ref{tab:heldout-results} reports qualified answers separately.

\begin{table}[htbp]
\centering
\begin{tabular}{lrrrrr}
\toprule
Reader & Complete & Conditional & Partial & Incorrect & Error\\
\midrule
Direct baseline & 44 & 1 & 4 & 11 & 0 \\
Source-checked & 36 & 11 & 4 & 6 & 3 \\
\bottomrule
\end{tabular}
\caption{Sixty held-out cases per condition. A conditional answer states an
assumption and gives the matching reference value. Such an answer does not
establish that the assumption is true. Reference alignment combines the first
two columns; it is distinct from unqualified correctness.}
\label{tab:heldout-results}
\end{table}

\begin{table}[htbp]
\centering
\begin{tabular}{lrr}
\toprule
Question type & Direct baseline & Source-checked\\
\midrule
Knowledge updates & 8/10 & 9/10 \\
Multi-session reasoning & 6/10 & 7/10 \\
Assistant recall & 10/10 & 8/10 \\
Personalized responses & 3/10 & 6/10 \\
User facts & 10/10 & 10/10 \\
Temporal reasoning & 8/10 & 7/10 \\
\bottomrule
\end{tabular}
\caption{Reference-aligned answers, including conditional answers, out of ten
cases per type. Partial answers and errors are excluded.}
\end{table}

The source-checked reader returned 47 reference-aligned answers, compared with
45 for the baseline. The two-case difference is too small to establish a
reliable advantage. Personalization improved, while several assistant-recall
and temporal cases deteriorated. The reader produced more conditional answers
and three validation errors. Two errors rejected inexact quotations; one
rejected an unresolved arithmetic endpoint. Counts of distinct events and
retrieval of every required event remain weaknesses.

The final source review also found a timezone presentation inconsistency: date
arithmetic used the configured timezone, while a displayed source timestamp
could use the database driver's timezone. The release normalizes the model's
anchors to the requested timezone. It also permits one inexact quotation to use
the existing repair allowance while preserving exact-match validation. A single
literal today or yesterday phrase can supply a conditional date candidate when
the extractor omitted an anchor; this does not establish the event date. These
post-measurement fixes are covered by regression checks and targeted replay;
they are not credited with changing the frozen sixty-case score. A further
release-model pass matched eleven of the twelve development references,
missing a location detail. The earlier 12/12 result was therefore not stable
across repeated model calls.

The larger natural-label study and the public study address different tasks.
Improved agreement with provisional retention labels does not establish improved
answer quality, and reference agreement does not certify truth outside the
supplied conversation. The learned classifiers remain advisory.

